\documentclass[dvipsnames]{article}

\usepackage[utf8]{inputenc}
\usepackage{fontspec}
\usepackage{polyglossia}
\usepackage{unicode-math}
\usepackage{csquotes}
\setdefaultlanguage{russian}
\setotherlanguage{english}

\setmainfont[SmallCapsFont=Latin Modern Roman Caps]{CMU Serif}
\setsansfont[SmallCapsFont=Latin Modern Roman Caps]{CMU Sans Serif}
\setmonofont[SmallCapsFont=Latin Modern Roman Caps]{CMU Typewriter Text}
% \setmathfont{Latin Modern Math}

\usepackage[letterpaper,top=2cm,bottom=2cm,left=3cm,right=3cm,marginparwidth=1.75cm]{geometry}

% Useful packages
\usepackage{amsmath}
\usepackage{graphicx}
\usepackage{amsthm}
\usepackage{array}
\usepackage{fancyhdr}
\usepackage{bussproofs}
\usepackage[colorlinks=true, allcolors=blue]{hyperref}
\usepackage[backend=biber]{biblatex}
\usepackage{xcolor}

\usepackage{minted2}
\usemintedstyle{xcode}
\setminted{escapeinside=||}
\newminted[stella]{cpp}{escapeinside=||,frame=single,texcomments,fontsize=\small}
\newmintinline[haskell]{haskell}{fontsize=\scriptsize}
\newmintinline[java]{java}{fontsize=\scriptsize}
\newmintinline[python]{python}{fontsize=\scriptsize}
\newmintinline[code]{c}{fontsize=\small}
\newmintinline[smallcode]{c}{fontsize=\small}
\let\smallcodeOrig\smallcode
\renewcommand{\smallcode}[2][]{%
  \ifmmode
    \if\relax\detokenize{#1}\relax
      \mathchoice{\mbox{\normalfont\small\ttfamily #2}}{\mbox{\normalfont\small\ttfamily #2}}{\mbox{\normalfont\small\ttfamily #2}}{\mbox{\normalfont\small\ttfamily #2}}%
    \else
      \mathchoice{\mbox{\normalfont\small\ttfamily #2}}{\mbox{\normalfont\small\ttfamily #2}}{\mbox{\normalfont\small\ttfamily #2}}{\mbox{\normalfont\small\ttfamily #2}}%
    \fi
  \else
    \smallcodeOrig[#1]{#2}%
  \fi}
\let\codeOrig\code
\renewcommand{\code}[2][]{%
  \ifmmode
    \if\relax\detokenize{#1}\relax
      \mathchoice{\mbox{\normalfont\small\ttfamily #2}}{\mbox{\normalfont\small\ttfamily #2}}{\mbox{\normalfont\small\ttfamily #2}}{\mbox{\normalfont\small\ttfamily #2}}%
    \else
      \mathchoice{\mbox{\normalfont\small\ttfamily #2}}{\mbox{\normalfont\small\ttfamily #2}}{\mbox{\normalfont\small\ttfamily #2}}{\mbox{\normalfont\small\ttfamily #2}}%
    \fi
  \else
    \codeOrig[#1]{#2}%
  \fi}

\usepackage{sourcecodepro}
\usepackage{mdframed}

\DeclareMathOperator{\xor}{\mathsf{xor}}
\DeclareMathOperator{\tru}{\mathsf{tru}}
\DeclareMathOperator{\fls}{\mathsf{fls}}

\addbibresource{../references.bib}

\begin{document}

\thispagestyle{fancy}
\lhead{\includegraphics[height=1cm]{../images/itmo_logo}}
\rhead{\includegraphics[height=1cm]{../images/stella-logo}}
\cfoot{Языки программирования. Семантика и системы типов (весна 2026, Университет ИТМО)}

\section*{Тема 11. Теоретическое задание \textcolor{blue}{(выполнил Павел Григоренко)}}%

\begin{enumerate}
  \item Предложите правила типизации на основе ограничений (см. \cite[\S 22.3]{TaPL-ru})
  для типов-произведений~\cite[\S 11.6]{TaPL-ru} и типов-сумм~\cite[\S 11.9]{TaPL-ru}.
  Затем, используя эти правила, выведите тип следующего терма в пустом контексте:

  \begin{center}
  \code{λy. λz. ((λx. case x of inl l ⇒ {l, y} | inr r ⇒ r) (inl z)).1}
  \end{center}

  А именно:
  \begin{enumerate}
    \item постройте дерево вывода типа с ограничениями для этого терма; в корне должно быть суждение вида $\cdot \vdash \code{t} : \code{S} \mid C$ для некоторых \code{S} и $C$;
    \item выпишите главный унификатор для множества ограничений $C$;
    \item выпишите главный тип терма \code{t}, используя полученный главный унификатор.
  \end{enumerate}

  \color{blue}
  \textbf{Решение.}
  \emph{место для решения}
  \color{black}

  \clearpage

  \item Рассмотрите следующие подстановки типовых переменных:
  \begin{align*}
    \xi &= [ \code{A} \mapsto \code{(B+Nat)→(C+Unit)},\; \code{C} \mapsto \code{(Nat→A)+Bool} ] \\
    \zeta &= [ \code{A} \mapsto \code{((Bool→Nat)+Nat)→(Nat+Unit)},\; \code{B} \mapsto \code{Bool→Nat},\; \code{C} \mapsto \code{(Nat→Unit)+Bool} ]
  \end{align*}

  Верно ли, что $\xi \sqsubseteq \zeta$? Верно ли, что $\zeta \sqsubseteq \xi$?
  Обоснуйте ваши ответы.

  \color{blue}
  \textbf{Решение.}
  \emph{место для решения}
  \color{black}

  \clearpage

  \item Покажите, что типы-суммы~\cite[\S 11.9]{TaPL-ru} можно представить в виде производной формы над Системой~F~\cite[\S 23.4]{TaPL-ru}:
  \begin{enumerate}
    \item покажите, как тип \code{A + B} можно заменить типом в Системе F (без расширений);
    \item выпишите раскрытие сокращений для термов \code{inl t} и \code{inr t}; при необходимости, используйте раскрытие сокращений для дерева вывода типа этого терма;
    \item выпишите раскрытие сокращений для терма \code{case t|$_1$| of inl x⇒t|$_2$| |$\mid$| inr x⇒t|$_3$|}; при необходимости, используйте раскрытие сокращений для дерева вывода типа;
    \item докажите, что раскрытие сокращений сохраняет семантику исходных правил типизации (\textsc{T-Inl}, \textsc{T-Inr} и \textsc{T-Case});
    \item докажите, что раскрытие сокращений сохраняет семантику исходных правил вычисления (\textsc{E-Inl}, \textsc{E-Inr}, \textsc{E-Case}, \textsc{E-CaseInl} и \textsc{E-CaseInr}).
  \end{enumerate}

  \color{blue}
  \textbf{Решение.}
  \emph{место для решения}
  \color{black}
\end{enumerate}

\section*{Список литературы}
\printbibliography[heading=none]

\end{document}
